\section{Hochschildの奇数次障害とrotation不変部分}\label{ho:sec:parity}
可換復元における対角smoothnessを、形式変形の障害と混同しないために、
滑らかな三次超曲面のKuznetsov成分の低次Hochschild群を全次元で整理する。
以下は既知のequivariant matrix factorizationの計算からの帰結であり、
残余dg圏の相対smoothnessに関する新たな一般定理ではない。

\begin{proposition}[奇数次Hochschild群とrotation固有値]\label{ho:prop:sectors}
$X=V(F)\subset\mathbb P^{N-1}_\CC$を滑らかな三次超曲面、$N=n+2\ge5$とし、
$J(F)=\CC[z_1,\dots,z_N]/(\partial F)$とする。
$\operatorname{Ku}(X)$の標準rotationを$\rho$とする。
$\zeta$を原始三乗根と書けば、次の複素ベクトル空間の分解が成り立つ：
\[
 \operatorname{HH}^{q}(\operatorname{Ku}(X))\cong
 \begin{cases}
 J(F)_{3q/2},&q\text{ 偶数},\\ 0,&q\text{ 奇数}
 \end{cases}
 \quad\oplus\quad
 \begin{cases}
 \CC_{\zeta}\oplus\CC_{\zeta^2},&3\mid N,\ q=N/3,\\
 0,&\text{それ以外}.
 \end{cases}
 \tag{HS}\label{ho:eq:sectors}
\]
ここで$J(F)_d=0$（$d<0$）とし、右側の二つの一次元空間は
$\rho$の共役作用の非自明な固有空間である。
$J(F)$から来る項では$\rho$は恒等的に作用する。
特に
\[
 \operatorname{HH}^3(\operatorname{Ku}(X))
 \cong\begin{cases}\CC^2&n=7,\\0&n\ne7,\end{cases}
 \qquad
 \boxed{\operatorname{HH}^3(\operatorname{Ku}(X))^\rho=0\quad(n\ge3).}
 \tag{V3}\label{ho:eq:HH3}
\]
さらに$\operatorname{HH}^2(\operatorname{Ku}(X))^\rho\cong J(F)_3$で、
$\dim_\CC J(F)_3=\binom N3$である。
\end{proposition}
\begin{proof}
Orlovのgraded matrix factorization表示と
Ballard--Favero--Katzarkov\cite[Theorem 5.39, pp.~82--85]{BFKER}の
固定sector公式を用いる。
$G=\mathbb G_m$は全生成変数の重みが$1$となる作用で、
potentialのcharacterは三次である。$K_\chi=\mu_3$とする。
$g=1$において$V_g=\mathbb A^N$、$W_g=0$である。
偏微分$(\partial F)$は原点にのみ共通零点を持つ長さ$N$の正則列であり、
Koszul cohomologyは零次の$J(F)$に集中する。
重み零の条件とcohomological gradingから、恒等sectorは偶数次数$q=2r$に
$J(F)_{3r}$を与え、奇数次数には寄与しない。

$g\in\mu_3\setminus\{1\}$なら固定空間は原点で、$\dim W_g=N$、
$J(F_g)=\CC$である。
Theorem 5.39のtwistとdeterminant characterを代入すると、
$u=(q-N)/2$に対する重み零条件は$3u+N=0$となる。
従って$q=N/3$でのみ各非自明sectorが一次元を与える。
各sectorは標準的な文字twist $(1)$ による共役作用で
$g$の文字（規約によって逆文字）を受ける。
これは同論文のkernelの固定graph分解と
\cite[Theorem 6.5]{BFKER}のcharacter作用の計算と同じものである。
Orlov同値の下で次数shift $(1)$ は$\rho$に対応するため、
二つのsectorの固有値は順序を除いて$\zeta,\zeta^2$となる。
$J(F)$から来る恒等sectorではこの作用は自明である。
これで\eqref{ho:eq:sectors}を得る。

$q=3$は奇数なので、非零になり得るのは$N/3=3$、すなわち$N=9$
（$n=7$）だけである。しかもその二方向には$\rho$-不変ベクトルがない。
$q=2$の不変部分は常に$J(F)_3$である。
滑らかさから偏微分は二次正則列なので、
$H_{J(F)}(t)=(1+t)^N$となり$\dim J(F)_3=\binom N3$を得る。
\end{proof}

\begin{corollary}[非楕円cubic fivefoldの変形障害の所在]\label{ho:cor:fivefold}
本研究の節点型・尖点型の二つの滑らかな可換特殊ファイバー
$X^5$について、共に
\[
 (\dim\operatorname{HH}^0,\dim\operatorname{HH}^1,
 \dim\operatorname{HH}^2,\dim\operatorname{HH}^3,
 \dim\operatorname{HH}^4)=(1,0,35,0,7).
 \tag{F5}\label{ho:eq:fivefold}
\]
従って滑らかな可換特殊ファイバーのdg圏自身の通常のMorita変形に対する
\emph{一次の}Kuranishi障害空間$\operatorname{HH}^3$は零である。
ただし、量子化から構成した$\Gamma_R$の
$\Gamma_R^e$上の対角perfect性を、これだけからは結論できない。
\end{corollary}
\begin{proof}
$N=7$を命題\ref{ho:prop:sectors}に代入する。
$J(F)_0,J(F)_3,J(F)_6$の次元はそれぞれ$1,35,7$で、
それ以外の表示次数$q=1,3$は零となる。
Morita変形の一次障害が$\operatorname{HH}^3$に入ることは
標準的なHochschild変形理論の適用である。
\end{proof}

\begin{remark}[残余対角との論理的境界]\label{ho:rem:barrier}
$\operatorname{HH}^3(\Ku(X))^\rho=0$は重要な形式的情報であるが、
この不変部分の消滅だけで任意のrotation付きdg圏変形の厳密な
$\mathbb Z/3$-equivariant obstruction theoryを構成したとはしない。
また、Artin切断上の変形が可換的に説明できたとしても、
有限型係数環上の\emph{実在する}残余対角のperfect性、
すなわち$\Gamma_R\in\Perf(\Gamma_R^{\mathrm{op}}\otL_R\Gamma_R)$は別の問題である。
所定の量子化のgeneric復元の最優先課題は依然としてこの残余対角のperfect性である。
本節はBFKの既知の公式の具体的な帰結と、その本研究における障害の切分けを与える。
\end{remark}
